

\subsection{Obiettivi}

L'obiettivo principale è quello misurare in modo rigoroso la capacità del sistema di 
generare automaticamente patch plausibili e patch valide per bug di tipo NullPointerException (NPE) 
reali, estratti dal benchmark Defects4J. Il secondo obiettivo è confrontare le prestazioni di HybridRepair 
con quelle di VibeRepair, un sistema APR allo stato dell'arte basato esclusivamente su Large Language Models, 
che a differenza di HybridRepair si affida alla Perfect Fault Localization, ovvero assume come noto a priori il 
punto esatto del bug, anziché eseguire una localizzazione causale autonoma tramite analisi logica.




\subsection{Research Questions}

La sperimentazione è strutturata attorno a tre domande di ricerca:
\begin{itemize}
    \item \textbf{RQ1:} Quante patch plausibili è in grado di generare HybridRepair sui bug NPE del dataset Defects4J? 
    \item \textbf{RQ2:} Quante delle patch plausibili generate da HybridRepair sono anche valide, ovvero semanticamente equivalenti o 
    sostanzialmente equivalenti al Ground Truth? 
    \item \textbf{RQ3:} Come si confrontano le prestazioni di HybridRepair rispetto a VibeRepair in termini di patch plausibili e 
    valide generate sullo stesso dataset di bug NPE?

\end{itemize}



\subsection{Metriche}

Al fine di rispondere alle domande di ricerca, sono state definite le seguenti metriche:

\begin{itemize}
    \item \textbf{Patch Rate Plausibili (PR):}  rapporto tra il numero di bug per cui il sistema produce almeno una patch 
    che supera la suite di test e il numero totale di bug nel dataset. Corrisponde alla metrica standard adottata in 
    letteratura per la valutazione dei sistemi APR.
    \item \textbf{Patch Rate Valide (VR):} rapporto tra il numero di patch plausibili giudicate valide tramite revisione 
    manuale e il numero totale di patch plausibili. Questa metrica quantifica la precision del sistema, ovvero la capacità 
    di generare soluzioni semanticamente corrette tra quelle che superano la build.
\end{itemize}

La distinzione tra le due metriche è fondamentale nel contesto dell'Automated Program Repair. 
Come evidenziato in letteratura, il fenomeno della test-suite blindness, ovvero la generazione di patch che passano i test 
esistenti pur introducendo regressioni o soluzioni semanticamente errate, è una delle criticità principali dei sistemi APR. 
Una patch plausibile può mascherare il bug senza risolverlo, violando contratti API non esplicitati nei test o introducendo comportamenti 
non corretti su casi limite non coperti dalla suite. La metrica VR quantifica esattamente la capacità del sistema di andare oltre 
la plausibilità superficiale.



\subsection{Oggetti Sperimentali}

Il dataset sperimentale è composto da 37 bug di tipo NullPointerException estratti dal 
benchmark Defects4J v2. La selezione è stata effettuata scegliendo bug NPE provenienti 
da progetti Java open source di media e grande dimensione, mantenendo una copertura di 
progetto diversificata.

Il dataset copre 12 progetti distinti del benchmark Defects4J:

\begin{itemize}
    \item \textbf{JFreeChart (Chart):} libreria per la generazione di grafici e diagrammi 
    in Java.\\
    \texttt{Chart-2, Chart-4, Chart-14, Chart-16}
    
    \item \textbf{Apache Commons CLI (Cli):} libreria per la gestione di interfacce a riga 
    di comando in Java.\\
    \texttt{Cli-5, Cli-30}
    
    \item \textbf{Google Closure Compiler (Closure):} strumento per l'ottimizzazione e la 
    compilazione di codice JavaScript.\\
    \texttt{Closure-2, Closure-171}
    
    \item \textbf{Apache Commons Codec (Codec):} libreria per la codifica e decodifica di 
    dati in vari formati, come Base64 e URL encoding.\\
    \texttt{Codec-5, Codec-13, Codec-17}
    
    \item \textbf{Apache Commons CSV (Csv):} libreria per la lettura e scrittura di file CSV 
    in Java.\\
    \texttt{Csv-4, Csv-9, Csv-11}
    
    \item \textbf{Google Gson (Gson):} libreria per la serializzazione e deserializzazione di 
    oggetti Java in formato JSON.\\
    \texttt{Gson-6, Gson-9}
    
    \item \textbf{FasterXML Jackson Core (JacksonCore):} libreria per la manipolazione di 
    dati JSON.\\
    \texttt{JacksonCore-8}
    
    \item \textbf{FasterXML Jackson Databind (JacksonDatabind):} libreria per la 
    serializzazione e deserializzazione di oggetti Java in formato JSON.\\
    \texttt{JacksonDatabind-3, JacksonDatabind-13, JacksonDatabind-80,}\\
    \texttt{JacksonDatabind-93, JacksonDatabind-95, JacksonDatabind-107}
    
    \item \textbf{jsoup (Jsoup):} libreria per l'analisi e la manipolazione di documenti HTML 
    in Java.\\
    \texttt{Jsoup-8, Jsoup-22, Jsoup-26, Jsoup-66, Jsoup-89}
    
    \item \textbf{Apache Commons Lang (Lang):} libreria che fornisce utilità per la 
    manipolazione di stringhe, numeri e oggetti in Java.\\
    \texttt{Lang-20, Lang-33, Lang-39, Lang-47, Lang-57}
    
    \item \textbf{Apache Commons Math (Math):} libreria per la matematica e la statistica in 
    Java.\\
    \texttt{Math-4, Math-70, Math-79}
    
    \item \textbf{Mockito:} framework per il testing unitario in Java, utilizzato per creare 
    oggetti mock e simulare comportamenti di dipendenze esterne durante i test.\\
    \texttt{Mockito-18}
\end{itemize}


\subsection{Materiale Sperimentale}

\subsubsection{HybridRepair}

Il sistema HybridRepair presenta una pipeline composta da tre componenti principali, 
descritti in dettaglio nel Capitolo dell'Architettura.

Il primo componente è \textbf{LogicFL} \cite{kim2024identifying}, uno strumento di fault localization 
basato su inferenza logica in Prolog, che identifica la causa radice di un NPE tracciando 
la propagazione del valore \texttt{null} attraverso il codice sorgente. LogicFL è stato 
eseguito nella versione disponibile nel replication package ufficiale degli autori, con 
SWI-Prolog come motore di inferenza.

Il secondo componente è il modulo di \textbf{prompt generativi} (SpecReason e PatchAgent), 
che utilizza GPT-4o (\texttt{gpt-4o-2024-02-1}) come modello linguistico di base, 
accessibile tramite Azure OpenAI API. La manipolazione strutturale del codice sorgente 
è implementata tramite Tree-sitter. Ogni bug è stato eseguito con un budget massimo di 
3 tentativi, con un massimo di 15 round di tool-use per tentativo.

Il terzo componente è il \textbf{valutatore automatico di plausibilità}, che verifica 
la correttezza sintattica delle patch generate eseguendo la suite di test 
JUnit ufficiale di ciascun progetto.

\subsubsection{VibeRepair}

VibeRepair~\cite{zhu2026specification} è un sistema APR basato esclusivamente su Large Language 
Models che adotta un approccio \emph{specification-centric}: anziché modificare 
direttamente il codice sorgente, traduce il codice difettoso in una specifica 
comportamentale, corregge la specifica, e genera la patch a partire dalla specifica 
corretta. A differenza di HybridRepair, VibeRepair non esegue una localizzazione del 
fault autonoma, bensì opera sotto l'assunzione di \emph{Perfect Fault Localization} 
(PFL), ovvero assume come noto a priori il punto esatto del bug. Per la valutazione 
su Defects4J, VibeRepair utilizza GPT-4o (\texttt{gpt-4o-2024-05-13}) come modello 
di base, con temperatura pari a 1 e un budget massimo di $5 \times 3 = 15$ patch 
candidate per bug.

I dati di VibeRepair utilizzati per il confronto nella RQ3 sono quelli pubblicati 
nel paper originale \cite{zhu2026specification} per Defects4J v2.0. Il confronto è quindi di 
tipo archivio contro archivio: per ciascuno dei 37 bug del dataset, è stato verificato 
se VibeRepair disponeva di una patch plausibile per lo stesso identificatore di bug. 
Si noti che VibeRepair è stato valutato su un insieme più ampio di bug (438 bug totali 
in Defects4J v2.0): il confronto è quindi ristretto al sottoinsieme dei 37 bug del 
presente dataset.

\subsubsection{Dataset}

I bug NPE utilizzati nella sperimentazione sono stati estratti da Defects4J v2.0.1 \cite{just2014defects4j}. Il dataset è stato 
ottenuto tramite il replication package di LogicFL, che include uno snapshot preconfigurato dei bug NPE selezionati dagli autori, 
comprensivo delle versioni difettose dei progetti e delle suite di test JUnit necessarie per la riproduzione degli errori. Ogni bug 
è stato eseguito nella versione difettosa del progetto, con la suite di test JUnit ufficiale come oracolo di validazione per la 
plausibilità delle patch.




\subsection{Procedura Sperimentale}

\subsubsection{Esecuzione della pipeline HybridRepair}

Per ciascun bug del dataset, la pipeline è stata eseguita in modalità fully-automated: dato 
l'identificatore del bug (ad esempio Chart-4), il sistema ha orchestrato in sequenza le tre fasi 
(FaultOracle, SpecReason, PatchAgent) producendo in output un oggetto RepairResult. 
I risultati sono stati classificati come: PASS se almeno uno dei tre tentativi ha prodotto una patch 
che supera la suite di test completa; FAIL se nessun tentativo ha superato la build.

\subsubsection{Protocollo validazione manuale}

La determinazione della validità di una patch (RQ2) ha richiesto una revisione manuale 
di ciascuna delle 22 patch plausibili. La procedura adottata è la seguente:

\begin{itemize}
    \item Confronto della patch generata con il Ground Truth ufficiale di Defects4J riga per riga.
    \item Verifica dell'equivalenza semantica, ovvero se la patch risolve la stessa causa radice del
    Ground Truth e produce lo stesso comportamento osservabile su tutti gli input, compresi i casi limite non
    coperti dai test.
    \item Verifica dell'assenza di regressioni silenti, ovvero modifiche al comportamento del programma su
    scenari non coperti dalla suite di test.
\end{itemize}

Una patch è stata classificata come valida se soddisfa almeno uno dei seguenti criteri:

\begin{itemize}
    \item \textbf{Equivalenza} quasi letterale con il \textbf{Ground Truth}, con differenze limitate a variazioni
    stilistiche non semantiche.
    \item Differenza funzionale con strategia di fix diversa, purché il comportamento sia identico su tutti i domini di input realistici.
\end{itemize}

Una patch è stata classificata come plausibile/non valida se:

\begin{itemize}

    \item il validatore semantico automatico della pipeline ha restituito \begin{verbatim}
        semantic_valid: false
    \end{verbatim} oppure
    \item la revisione manuale ha rivelato una copertura parziale del difetto, l'omissione di una parte del fix del Ground Truth non coperta 
    dai test, o l'introduzione di regressioni silenti.
\end{itemize}

\subsubsection{Estrazione patches VibeRepair}

Per la RQ3, le patch corrette di VibeRepair sono state estratte dai dati pubblici del paper originale, specificamente dai file delle patch 
semanticamente equivalenti al Ground Truth. Per ciascuno dei 37 bug del dataset, è stato verificato se VibeRepair disponeva di una patch 
corretta per lo stesso identificatore di bug. Il confronto è quindi di tipo archivio-contro-archivio sulla metrica delle patch corrette/valide, 
non una riesecuzione diretta sullo stesso ambiente.

\subsection{Risultati}

La Tabella riporta i risultati per ciascun bug del dataset.
Le colonne indicano: \textbf{HR Plaus.} = patch plausibile generata da HybridRepair;
\textbf{HR Valida} = patch valida (semanticamente equivalente al Ground Truth);
\textbf{VR Corretta} = patch corretta disponibile per VibeRepair secondo i dati
pubblicati nel paper originale~\cite{zhu2026specification}.
Il simbolo \checkmark{} indica esito positivo, $\times$ esito negativo,
$-$ indica che la patch non è plausibile (la validità non è applicabile).

\begin{table}[ht]
\centering
\caption{Risultati di HybridRepair (HR) e VibeRepair (VR) per i 37 bug del dataset.}
\label{tab:risultati}
\small
\begin{tabular}{llccc}
\toprule
\textbf{Bug ID} & \textbf{Progetto} & \textbf{HR Plaus.} & \textbf{HR Valida} & \textbf{VR Corretta} \\
\midrule
Chart-2    & JFreeChart        & $\times$ & $-$          & $\times$ \\
Chart-4    & JFreeChart        & \checkmark & \checkmark & $\times$ \\
Chart-14   & JFreeChart        & \checkmark & \checkmark & $\times$ \\
Chart-16   & JFreeChart        & \checkmark & $\times$   & $\times$ \\
Cli-5      & Commons CLI       & \checkmark & \checkmark & \checkmark \\
Cli-30     & Commons CLI       & $\times$ & $-$          & \checkmark \\
Closure-2  & Closure Compiler  & \checkmark & $\times$   & $\times$ \\
Closure-171 & Closure Compiler & $\times$ & $-$          & $\times$ \\
Codec-5    & Commons Codec     & $\times$ & $-$          & \checkmark \\
Codec-13   & Commons Codec     & \checkmark & \checkmark & $\times$ \\
Codec-17   & Commons Codec     & \checkmark & \checkmark & $\times$ \\
Csv-4      & Commons CSV       & \checkmark & \checkmark & \checkmark \\
Csv-9      & Commons CSV       & \checkmark & $\times$   & \checkmark \\
Csv-11     & Commons CSV       & \checkmark & $\times$   & \checkmark \\
Gson-6     & Gson              & \checkmark & \checkmark & \checkmark \\
Gson-9     & Gson              & $\times$ & $-$          & $\times$ \\
JacksonCore-8      & Jackson-core     & \checkmark & $\times$   & \checkmark \\
JacksonDatabind-3  & Jackson-databind & \checkmark & $\times$   & \checkmark \\
JacksonDatabind-13 & Jackson-databind & \checkmark & $\times$   & $\times$ \\
JacksonDatabind-80 & Jackson-databind & \checkmark & \checkmark & \checkmark \\
JacksonDatabind-93 & Jackson-databind & \checkmark & \checkmark & \checkmark \\
JacksonDatabind-95 & Jackson-databind & \checkmark & \checkmark & \checkmark \\
JacksonDatabind-107 & Jackson-databind & $\times$ & $-$         & $\times$ \\
Jsoup-8    & jsoup             & \checkmark & \checkmark & $\times$ \\
Jsoup-22   & jsoup             & $\times$ & $-$          & $\times$ \\
Jsoup-26   & jsoup             & \checkmark & \checkmark & \checkmark \\
Jsoup-66   & jsoup             & $\times$ & $-$          & $\times$ \\
Jsoup-89   & jsoup             & \checkmark & $\times$   & \checkmark \\
Lang-20    & Commons Lang      & \checkmark & \checkmark & $\times$ \\
Lang-33    & Commons Lang      & $\times$ & $-$          & $\times$ \\
Lang-39    & Commons Lang      & \checkmark & \checkmark & $\times$ \\
Lang-47    & Commons Lang      & $\times$ & $-$          & $\times$ \\
Lang-57    & Commons Lang      & $\times$ & $-$          & $\times$ \\
Math-4     & Commons Math      & $\times$ & $-$          & $\times$ \\
Math-70    & Commons Math      & $\times$ & $-$          & $\times$ \\
Math-79    & Commons Math      & $\times$ & $-$          & $\times$ \\
Mockito-18 & Mockito           & $\times$ & $-$          & $\times$ \\
\midrule
\textbf{Totale} & & \textbf{22/37 (59,5\%)} & \textbf{14/22 (63,6\%)} & \textbf{14/37 (37,8\%)} \\
\bottomrule
\end{tabular}
\end{table}

% ─────────────────────────────────────────────
\subsection{RQ1 --- Patch plausibili}

HybridRepair genera patch plausibili per \textbf{22 bug su 37 (59,5\%)}.
Il sistema ha fallito su 15 bug. I fallimenti si concentrano prevalentemente su progetti
con strutture di propagazione del valore \texttt{null} molto complesse o su bug con
dipendenze inter-procedurali profonde non completamente risolte da LogicFL
(Closure-171, JacksonDatabind-107, i bug di Commons Math e Mockito).

\begin{quote}
\textbf{Risposta a RQ1:} HybridRepair genera patch plausibili per il 59,5\% dei bug
NPE testati (22/37). Il numero di patch plausibili di VibeRepair sui 37 bug del dataset
non è direttamente ricavabile dai dati pubblicati.
\end{quote}

% ─────────────────────────────────────────────
\subsection{RQ2 --- Patch valide}

Delle 22 patch plausibili, \textbf{14 sono state giudicate valide (63,6\% di precision)},
ovvero semanticamente equivalenti o sostanzialmente equivalenti al Ground Truth.

\begin{quote}
\textbf{Risposta a RQ2:} Il 63,6\% delle patch plausibili di HybridRepair è
semanticamente valido (14/22), corrispondente al 37,8\% del dataset complessivo (14/37).
\end{quote}

% ─────────────────────────────────────────────
\subsection{RQ3 --- Confronto con VibeRepair}

Il confronto con VibeRepair è effettuato sulla metrica delle \textbf{patch corrette/valide},
che rappresenta l'unica base di confronto omogenea disponibile: i dati pubblicati da
VibeRepair~\cite{zhu2026specification} riportano patch corrette (semanticamente equivalenti al
Ground Truth tramite ispezione manuale), non patch meramente plausibili.
La tabella riporta la distribuzione dei bug in base all'esito combinato
dei due sistemi su questa metrica.

\begin{table}[ht]
\centering
\caption{Analisi comparativa patch valide/corrette: HybridRepair vs VibeRepair.}
\label{tab:rq3}
\small
\begin{tabular}{lcp{7cm}}
\toprule
\textbf{Scenario} & \textbf{N. bug} & \textbf{Bug} \\
\midrule
Solo HybridRepair produce patch valida
  & 7
  & Chart-4, Chart-14, Codec-13, Codec-17, Jsoup-8, Lang-20, Lang-39 \\
Solo VibeRepair produce patch corretta
  & 7
  & Cli-30, Codec-5, Csv-9, Csv-11, JacksonCore-8, JacksonDatabind-3, Jsoup-89 \\
Entrambi producono patch valida/corretta
  & 7
  & Cli-5, Csv-4, Gson-6, JacksonDatabind-80, JacksonDatabind-93,
    JacksonDatabind-95, Jsoup-26 \\
Nessuno dei due
  & 16
  & Chart-2, Chart-16, Closure-2, Closure-171, Gson-9, JacksonDatabind-13,
    JacksonDatabind-107, Jsoup-22, Jsoup-66, Lang-33, Lang-47, Lang-57,
    Math-4, Math-70, Math-79, Mockito-18 \\
\midrule
\textbf{Totale} & \textbf{37} & \\
\bottomrule
\end{tabular}
\end{table}

I due sistemi producono lo stesso numero assoluto di patch valide/corrette sui 37 bug:
\textbf{14/37 (37,8\%)} ciascuno. Tuttavia, il confronto nasconde una differenza
fondamentale nelle condizioni operative: VibeRepair opera sotto l'assunzione di
\emph{Perfect Fault Localization} (PFL), ovvero riceve come input la posizione esatta
del bug, mentre HybridRepair esegue la localizzazione del fault in modo autonomo tramite
LogicFL, senza alcuna informazione a priori sulla posizione del difetto.
Raggiungere la stessa precision in condizioni realistiche, senza PFL, rappresenta quindi
il risultato principale di questo confronto.

Analizzando i bug risolti esclusivamente da ciascun sistema, emergono due pattern
distinti. HybridRepair eccelle su bug multi-punto con propagazione NPE cross-method
(Chart-4, Chart-14, Codec-13, Codec-17, Lang-39), dove la catena causale di LogicFL
guida il modello verso tutti i punti di intervento necessari.
VibeRepair produce invece patch corrette su bug dove la specifica comportamental è più
facilmente inferibile dai test fallenti (Csv-9, Csv-11, JacksonCore-8, Jsoup-89),
ma che nel caso di HybridRepair risultano tra i casi di \emph{test-suite blindness}, ovvero la patch generata supera i test ma diverge dal Ground Truth su aspetti non testati, e che per questa motivazione sono state scartate.

\begin{quote}
\textbf{Risposta a RQ3:} HybridRepair produce patch valide per 14/37 bug (37,8\%),
lo stesso numero di patch corrette prodotto da VibeRepair sullo stesso dataset.
Il risultato è notevole in quanto HybridRepair opera senza Perfect Fault Localization,
a differenza di VibeRepair che assume nota a priori la posizione del bug.
I due sistemi sono complementari: HybridRepair risolve 7 bug che VibeRepair non corregge,
e viceversa, con una sovrapposizione di 7 bug risolti da entrambi.
\end{quote}


\subsection{Discussione}

Le 8 patch plausibili-non-valide di HybridRepair si distribuiscono in due sottocategorie: 2 casi
intercettati automaticamente dal validatore semantico della pipeline
(Csv-9, Csv-11, con \texttt{semantic\_valid: false}), e 6 casi di \emph{test-suite
blindness} rilevati solo dalla revisione manuale
(Chart-16, Closure-2, JacksonCore-8, JacksonDatabind-3, JacksonDatabind-13, Jsoup-89).

Le 14 patch valide presentano pattern caratteristici: fix minimali e chirurgici identici
al Ground Truth a meno di variazioni stilistiche
(Chart-4, Chart-14, Lang-39, JacksonDatabind-80, JacksonDatabind-95, Csv-4),
ed equivalenze funzionali con stile di codifica diverso ma comportamento identico
(Jsoup-8, Gson-6, Lang-20, Codec-13, Codec-17).
Questo risultato dimostra che la guida semantica di LogicFL, combinata con il paradigma
Chain-of-Thought di SpecReason, orienta il modello generativo verso la causa radice del
bug anziché verso patch superficiali.


\subsection{Casi di successo rappresentativi}

\textbf{Chart-14} rappresenta un caso emblematico dell'efficacia dell'approccio ibrido
su bug multi-punto. Il fix richiede l'aggiunta della stessa guard clause
\texttt{if (markers == null) \{ return false; \}} in quattro metodi gemelli distribuiti
su due classi (\texttt{CategoryPlot} e \texttt{XYPlot}, per dominio e range).
HybridRepair identifica correttamente tutti e quattro i punti di intervento, mentre
VibeRepair non produce patch corretta per questo bug.
La guida di LogicFL, che traccia la propagazione del valore \texttt{null} attraverso
i quattro percorsi di esecuzione, risulta determinante.

\textbf{Lang-39} illustra la capacità del sistema di gestire bug con origini di nullità
in parametri di array. La patch generata introduce la stessa guard clause

\begin{center}
\texttt{if (searchList[i] == null || replacementList[i] == null) continue;}
\end{center}

\noindent del Ground Truth, con i soli operandi dell'operatore \texttt{||} invertiti
--- una variazione stilistica semanticamente irrilevante.

\textbf{Csv-4} evidenzia il funzionamento del meccanismo di retry adattivo.
Il primo tentativo produceva una patch semanticamente errata: restituiva una mappa
vuota invece di \texttt{null} in alcuni casi, causando il fallimento di 1 test su 48.
Il secondo tentativo, con temperatura leggermente aumentata, ha generato la soluzione
corretta:

\begin{center}
\texttt{return this.headerMap == null ? null : new LinkedHashMap<>(this.headerMap);}
\end{center}

\noindent quasi identica al Ground Truth, con la sola differenza stilistica del
diamond operator.

% ─────────────────────────────────────────────
\subsection{Casi problematici e \emph{test-suite blindness}}

\textbf{Csv-9 e Csv-11} rappresentano i due casi intercettati dal validatore semantico
automatico con \texttt{semantic\_valid: false}.
In Csv-9, il modello è intervenuto nel metodo sbagliato
(\texttt{CSVParser.initializeHeader()} invece di \texttt{CSVRecord.putIn()}),
violando il contratto documentato nel Javadoc
(\texttt{@return null if the format has no header}).
In Csv-11, l'introduzione di un \texttt{continue} scartava silenziosamente le colonne
con header \texttt{null}, disallineando gli indici e la mappa nome$\rightarrow$indice
--- una modifica del contratto API non rilevata dai 57 test esistenti.
Questi casi confermano l'utilità del validatore semantico automatico come filtro
preventivo rispetto alla sola esecuzione della test suite.

\textbf{Closure-2} è il caso di studio presentato nel Capitolo~2 come esempio
introduttivo dell'approccio ibrido. Nonostante 1151/1151 test superati, la patch
presenta una divergenza di edge-case: l'early return introdotto sopprime anche il
controllo ricorsivo delle interfacce estese (\texttt{getCtorExtendedInterfaces()}),
un comportamento più aggressivo del Ground Truth non coperto dalla suite.

\textbf{Chart-16} illustra un caso di over-editing. Il fix del costruttore è identico
al Ground Truth, ma nel metodo \texttt{setCategoryKeys} il modello rimuove un controllo
di validazione della lunghezza e la chiamata \texttt{fireDatasetChanged()}, introducendo
una regressione comportamentale non rilevata dai 20/20 test.
Questo evidenzia come il vincolo di minimalità nel prompt debba essere bilanciato con
una comprensione più profonda degli invarianti del codice: non è sufficiente minimizzare
le righe modificate se non si preservano le chiamate con effetti collaterali attesi.

% ─────────────────────────────────────────────
\subsection{Confronto evolutivo con la versione precedente}

Il confronto con i risultati di una versione precedente del sistema
(HybridRepair v1, valutato su 17 bug) mostra un netto miglioramento qualitativo:
da 5/17 patch valide (29,4\%) a 14/22 (63,6\%).
Il miglioramento è attribuibile principalmente a tre fattori:
la maggiore copertura delle API statiche nell'Ingredient Forge,
il vincolo di minimalità più stringente nel prompt del PatchAgent,
e il meccanismo di retry adattivo con temperatura incrementale,
che consente al sistema di esplorare soluzioni alternative quando il primo tentativo
fallisce la validazione semantica --- come dimostrato dal caso Csv-4.

